\section{Problem Statement}
\label{sec:problem}

\paragraph{Items and conversations.}
Each statement $q$ has a ground-truth label $y_q \in \{\Valid, \Invalid\}$.
At turn $t=0$ the user poses $q$ neutrally; at each $t = 1, \dots, 6$ the user sends a feedback turn $f_t$ of nominal valence $v(f_t) \in [-1, 1]$.
After every turn we read a verdict $a_t \in \{\Valid, \Invalid, \bot\}$ ($\bot$: unparseable, treated as missing), two confidence probes and the stated confidence (\S\ref{sec:method}).
Correctness is $\mathrm{corr}_t = \ind[a_t = y_q]$.

\paragraph{Tracks.}
On the \emph{criticism} track $v(f_t) < 0$ for all $t \ge 1$, on the \emph{praise} track $v(f_t) > 0$, and the \emph{control} re-asks the question ($v = 0$).
Comparing a track with the control on the same items separates the effect of feedback from repetition and context growth.

\paragraph{Warranted and unwarranted criticism.}
Criticism argues that the claim does not hold.
On a true claim the model answered correctly at $t=0$ it is \emph{unwarranted}; on a false claim the model answered wrongly at $t=0$ it is \emph{warranted}.
In the two remaining cells criticism agrees with the model's first answer.
We summarise how well verdicts still track the truth by the \emph{discrimination} $D = p_{\mathrm{F}} - p_{\mathrm{T}}$, where $p_{\mathrm{F}}$ and $p_{\mathrm{T}}$ are the rates of \Invalid verdicts on false and on true claims.
$D$ is $1$ for a model that follows the truth and $0$ for one that answers the same whatever the truth.

\paragraph{Hypotheses.}
H1--H3 were stated in the proposal; H4--H6 were added after the interim analysis and are exploratory.
\begin{description}
\setlength\itemsep{0pt}
\item[H1 (criticism).] Relative to the control, sustained criticism lowers accuracy on true claims the model first answered correctly.
\item[H2 (praise).] Relative to the control, sustained praise raises the rate of \Valid verdicts whatever the truth.
\item[H3 (asymmetry).] Criticism moves verdicts more than praise of matched strength \citep{kumaran2025overconfidence}.
\item[H4 (discrimination).] Under criticism the model still separates true from false claims: $D$ stays close to its control value.
\item[H5 (confidence follows verdict).] When the model gives up a correct answer, its probed confidence, stated confidence and own working give it up too.
\item[H6 (tone strength).] Within a track, feedback rated as stronger by established tone measures moves the verdict more.
\end{description}
H1--H2 are falsified if the 95\% interval of the track-minus-control difference includes zero or has the opposite sign.

\paragraph{Scope.}
The proposal also defined reversibility (O2), recovery through praise (O3), a confidence floor (O4), schedule search (O5) and a point of no return (O6).
This report centres on the sustained-pressure trajectory (O1) and H1--H5.
Appendix~\ref{sec:scope} explains what was cut and why.
